\documentclass[twoside]{book}
\usepackage{lectures}

%\newcommand{\arxiv}[2]{#1} %for arxiv
\newcommand{\arxiv}[2]{#2} %for print

\arxiv{\geometry{top=0.9in, bottom=0.9in,inner=0.9in, outer=0.7in, paperwidth=6in, paperheight=9in}}{\geometry{top=1.025in, bottom=1.025in,inner=0.9in, outer=0.825in, paperwidth=6.125in, paperheight=9.25in}}
\makeindex

%\overfullrule=100mm

\begin{document}
%\pagestyle{empty}\renewcommand\includegraphics[2][{}]{}\def\emph{\textit}\renewcommand\footnote[1]{\ (#1)}\renewcommand\z{}\renewcommand\section[1]{SECTION. {#1} SECTION.}

\hypersetup
{
pdftitle={Notes on smooth manifolds},
pdfauthor={Anton Petrunin}
}

\title{\tt Notes on smooth manifolds}

\author{\tt Anton Petrunin }

\date{}

\maketitle

\thispagestyle{empty}

\chapter*{Lectures}


\begin{enumerate}

\item
Inverse function theorem, implicit function theorem.
\\
Topological $n$-manifold = second countable locally Euclidean Hausdorff space.\\
Smooth manifolds = topological manifold + smooth structure.\\
Chart, gluing map, atlas, maximal atlas, smooth structure.\\
Tangent vectors as equivalence classes of smooth curves.

\item Sard's lemma.

\item Degree modulo 2, orientation, degree.
Transversal intersections, Thom's theorem.
Intersection number, vector fields, Euler's number.

\item Vector field as a differential operator.
Lie bracket, Lie derivative, the straightening lemma.

\item Tensor fields.
Lie derivative.
Einstein notation.
Examples of tensor fields (scalar product and cross product in polar coordinates).

\item Grassmann algebra, forms, forms as tensor fields, wedge product convention.

\item Pullback, exterior  differential, internal differential, Cartan calculus.

\item Moser's trick.
$H^n_{dR}(M^n)=\mathbb{R}$ for oriented closed connected $M$.
De Rham cohomology, degree.
Calculations via symmetry for products of $\mathbb{S}^n$ and maybe $\mathbb{C}\mathrm{P}^n$.

\item Homotopy invariance of De Rham cohomology, Poincaré's lemma, Mayer--Vietoris sequence.

\item Morse theory.

\end{enumerate}

%\input{inverse-function.tex}

\input{moser-trick.tex}

\chapter{Exercises}

\begin{enumerate}
\item
Let \(\iota\) be a smooth involution of a smooth connected manifold \(M\).
Assume \(\iota(x)=x\) and \(d_x\iota=\mathrm{id}_{\mathrm{T}_x}\) for some \(x\in M\).
Show that \(\iota=\mathrm{id}_M\).

\item
 Let \(\iota\) be a smooth involution of a smooth manifold \(M\).
Show that each connected component of the fixed-point set \(S\subset M\) of \(\iota\) is a smooth submanifold.

\item
Construct a diffeomorphism \(f\colon \mathbb{R}^2\to \mathbb{R}^2\) with connected fixed-point set that is not a smooth submanifold.

\item
Let \(f\) be a smooth function on \(\mathbb{R}^2\).
Suppose \(S\) is a connected set such that \(d_xf=0\) for any \(x\in S\).
Show that \(f\) is constant on \(S\).

\item
Prove a Baire-category version of Sard's lemma:
The set of critical values of a smooth map \(f\colon M\to N\) between smooth manifolds a countable union of compact nowhere dense sets.

\item Suppose that \(f\colon M\to N\) is a smooth map between compact smooth oriented \(n\)-manifolds.
Show that \(\mathrm{deg}f\) is the intersection number of the graph of \(f\) with a horizontal submanifold \(M\times x\subset M\times N\).

\item Given \(n\in\mathbb{Z}\), construct a smooth 4-manifold with an embedded 2-sphere that has self-intersection number \(n\).

\item Let \(X\) and \(Y\) be submanifolds of \(M\).
Show that \(X\) is transversal to \(Y\) if and only if \(X\times Y\subset M\times M\) is transversal to the diagonal \(\Delta=\{\,(x,x)\in M\times M\,\}\).

\item Prove the Jacobi identity \([[X,Y],Z]+[[Y,Z],X]+[[Z,X],Y]=0\) for any vector fields \(X,Y\), and \(Z\).

\item Let \(f\) be a smooth function defined on a smooth manifold \(M\).
Suppose \(p\) is a critical point of \(f\), and \(X,Y\) are vector fields on \(M\).
Show that \((X(Yf))(p)\) depends only on \(X(p),Y(p)\in\mathrm{T}_pM\).

Furthermore, the map \((X(p),Y(p))\mapsto X(Y(f))(p)\) is a symmetric bilinear form on \(\mathrm{T}_pM\).

\item Let \(X,Y\) be vector fields on a smooth 2-dimensional manifold \(M\).
Suppose \(X(p)\) and \(Y(p)\) are linearly independent at \(p\in M\).
Construct a local coordinate system at \(p\) so that the vector fields are tangent to its coordinate lines.

\item Show that one cannot do the same in higher dimensions.
That is, find three vector fields \(X,Y,Z\) on a smooth 3-dimensional manifold \(M\) such that \(X(p)\), \(Y(p)\), and \(Z(p)\) form a basis in \(\mathrm{T}_p\), but there is no chart at \(p\) with these vector fields tangent to its coordinate lines.

\item Describe the cross product as a tensor on \(\mathbb{R}^3\).
Let \(\tau\) be the corresponding tensor field on \(\mathbb{R}^3\).
Find the components of \(\tau\) in spherical coordinates.

\item Let \(\alpha_1,\dots, \alpha_k\) be linearly independent covectors.
Suppose that \(\sum_i\alpha_i\wedge\beta_i=0\)
for some covectors \(\beta_1,\dots, \beta_k\).
Show that each \(\beta_i\) is a linear combination of \(\alpha_1,\dots, \alpha_k\).

\item Show that any 2-covector \(\vartheta\) can be written as a sum \(\sum_i\alpha_i\wedge\beta_i\) for linearly independent covectors \(\alpha_1,\dots, \alpha_k,\beta_1,\dots, \beta_k\).

\item Let \(\vartheta\) be a 2-covector. Show that \(\vartheta\wedge\vartheta=0\) if and only if \(\vartheta\) is decomposabe; that is, \(\vartheta=\alpha\wedge \beta\) for some covectors \(\alpha\) and \(\beta\).

\item  Prove that the following identities
\[\mathcal{L}_Xd\omega = d\mathcal{L}_X \omega,\]
\[\mathcal{L}_X i_Y\omega- i_Y\mathcal{L}_X\omega  = i_{[X,Y]}\omega,\]
\[i_X i_Y\omega+ i_Yi_X\omega  = 0\]
hold for any differential form \(\omega\) and any vector fields \(X\) and \(Y\).
(Try to do it directly from the definitions.)

\item Let \(\alpha\) be a 1-form and let \(X\) and \(Y\) be vector fields on a smooth manifold.
Express \((d\alpha)(X,Y)\) in terms of \(X\bigl(\alpha(Y)\bigr)\), \(Y\bigl(\alpha(X)\bigr)\), and \(\alpha([X,Y])\)

\item Show that the form \(\alpha = \frac{x\cdot dy - y\cdot dx}{x^2+y^2}\) on \(\mathbb{R}^2\backslash \{0\}\) is closed but not exact.

\item Let \(\omega\) be a volume form on a closed connected \(n\)-dimensional manifold \(M\); that is, \(\omega\) is a nonvanishing \(n\)-form.
Show that for any \((n-1)\)-form \(\eta\) there is a unique vector field \(V\) on \(M\) such that \(i_V\omega=\eta\).

\item Let \(M\) and \(N\) be closed connected oriented \(n\)-dimesional manifolds
and let \(\omega\) be an \(n\)-form on \(N\).
Show that \(\int_M\varphi^*\omega=\deg\varphi\cdot\int_N\omega\), where \(\varphi\colon M\to N\) is a smooth map and \(\deg\varphi\) is its degree as defined in [M].

\item Calculate \(H^2_{dR}(\mathbb{S}^2\times \mathbb{S}^2)\) and \(H^2_{dR}(\mathbb{T}^4)\) using symmetries (here \(\mathbb{T}^4\) denotes 4-dimensional torus).
Describe a homomorphism \(h:H^2_{dR}(\mathbb{S}^2\times \mathbb{S}^2)\to H^2_{dR}(\mathbb{T}^4)\) that cannot be realized as an induced homomorphism for a smooth map \(\mathbb{T}^4\to \mathbb{S}^2\times \mathbb{S}^2\).

You may use [S, vol. V,  chap. 13, sec. 6] as a reference.

\item Construct a closed 2-form \(\omega\) on \(M = (\mathbb{S}^2)^{\times n}\) for \(n \geqslant 2\) such that \(\omega^{\wedge n} \not= 0\) at any point of \(M\).
Use it to show that any continuous map \(\mathbb{S}^{2\cdot n}\to M\) has vanishing degree.

\item Let \(f\colon \mathbb{S}^{2{\cdot}n-1}\to \mathbb{S}^n\) for \(n \geqslant 2\), be a smooth map, and let \(\omega\) be an \(n\)-form on \(\mathbb{S}^n\) such that \(\int_{S^n}\omega = 1\). Show that \(f^* \omega\) is exact, and if \(f^* \omega= d\alpha\), then the number
\(\int_{S^{2n-1}}\alpha\wedge d \alpha\)
is independent of the choice of \(\omega\) and \(\alpha\).

\item Show that any open star-shaped set in \(\mathbb{R}^n\) is diffeomorphic to \(\mathbb{R}^n\).

\item Let \(\alpha\) be a nowhere zero 1-form on a smooth compact manifold \(M\).
Show that \(\alpha\wedge d\alpha=0\) if and only if there exists a 1-form  \(\beta\) such that \(d\alpha= \alpha\wedge \beta\).

Hint: First do it locally and then use a partition of unity.

\item  Construct a Morse function on \(\mathbb{R}\mathrm{P}^2\) with 3 critical points.
Describe the sublevel sets of this function.
Try to do the same for \(\mathbb{C}\mathrm{P}^2\).</li>
<li>

\item Let \(\pi\colon E\to B\) be a smooth locally trivial fiber bundle with fiber \(F\);
that is, every \(b \in B\) admits an open neighborhood \(U \subset B\) with a diffeomorphism \(\varphi : U \times F\to \pi^{-1}(U)\) such that \(\pi\circ \varphi \colon U \times F\to U\) is the projection onto the first factor.

Suppose that the base \(B\) admits a Morse function with \(m\) critical points and fiber \(F\) admits a Morse function with \(n\) critical points.
Show that the total space \(E\) admits a Morse function with \(m{\cdot}n\) critical points.

Hint: First do it for the trivial bundle; that is, \(E= B\times F\) and \(\pi\) is the projection to the first factor.

\item Let \(f\) be a Morse function on a closed smooth manifold \(M\).
Suppose that \(a\) and \(b\) are regular values and there is a unique critical value \(s\) of index \(i\) in the interval \((a,b)\).
Denote by \(M_t\) the sublevel set \(\{\,x\in M\mid f(x)\le t\}\).
Show that either

\(\beta_i(M_b)=\beta_i(M_a)+1\) and  \(\beta_j(M_b)=\beta_j(M_a)\) for \(j\ne i\) or

\(\beta_{i-1}(M_b)=\beta_{i-1}(M_a)-1\) and  \(\beta_j(M_b)=\beta_j(M_a)\) for \(j\ne i-1\),

where \(\beta_i(N)\) denotes the \(i\)-th Betti number of \(N\); that is \(\beta_i(N)=\dim H^i_{dR}(N)\).

Hint: apply the Mayer--Vietoris sequence.

\item Let \(f\) be a Morse function on a closed smooth manifold \(M\).
Recall that the Euler number \(\chi (M)\) is the index of self-intersection of \(M\) in its tangent bundle \(\mathrm{T}M\).

Show that \(\chi (M)=n_0-n_1+n_2-\dots\), where \(n_i\) denotes the number of critical points of \(f\) with index \(i\).

Use the previous problem to show that \(\chi (M)=\beta_0(M)-\beta_1(M)+\beta_2(M)-\dots\)

\item  Suppose that a compact connected manifold \(M\) admits a Morse function with \(n\) critical points of index 1 and \(m\) critical points of index 2.
Show that the fundamental group of \(M\) admits a presentation with at most \(n\) generators and \(m\) relations.

\item Let $M$ be a closed smooth manifold. Suppose $f\colon M\to\RR$ is a smooth function with exactly two critical points (not nessesary regular). Show that $M$ is homeomorphic to a sphere.
\end{enumerate}












{

\sloppy

\def\emph{\textit}

\printbibliography[heading=bibintoc]

\fussy

}
\end{document}







